\title{DeepSeekMath: Pushing the Limits of Mathematical Reasoning in Open Language Models}

\begin{abstract}
The inherently intricate and highly structured demands of mathematical reasoning remain a formidable obstacle for language models. This study presents DeepSeekMath~7B, an extension of DeepSeek-Coder-Base-v1.5 7B achieved through continued pre-training with 120B math-specific tokens derived from Common Crawl, incorporating both natural language and programming code. DeepSeekMath~7B attains a notable 51.7\% accuracy on the advanced MATH benchmark—achieved independently of external toolkits or ensemble (voting) methods—placing its performance within close range of leading models like Gemini-Ultra and GPT-4. When applying self-consistency over 64 generations, DeepSeekMath~7B reaches an even higher 60.9\% on MATH. The superior mathematical reasoning exhibited by DeepSeekMath~originates from two primary innovations: firstly, the extensive leveraging of publicly accessible web data using a rigorously designed data curation process, and secondly, the implementation of Group Relative Policy Optimization (GRPO), a modified approach to Proximal Policy Optimization (PPO) that advances both mathematical reasoning and the efficiency of PPO's memory management.
\end{abstract}

\section{Introduction}

Large language models (LLM) have transformed methodologies for mathematical reasoning within artificial intelligence, leading to major progress on benchmarks such as quantitative reasoning \citep{MATH} and geometry reasoning \citep{trinh2024solving}. Additionally, these models have played a key role in enabling humans to tackle intricate mathematical challenges \citep{tao}. Despite these achievements, state-of-the-art systems like GPT-4 \citep{gpt4} and Gemini-Ultra \citep{gemini} remain proprietary, and the open-source alternatives currently available exhibit much lower performance levels.

This paper presents DeepSeekMath, a specialized language model for mathematics that markedly surpasses the mathematical reasoning capabilities of other open-source models and closely rivals GPT-4 on several academic benchmark tasks. Central to this achievement is the construction of the DeepSeekMath Corpus—a comprehensive, high-caliber pre-training resource encompassing 120B mathematical tokens. The data is sourced from Common Crawl (CC) by deploying a fastText-based classifier \citep{joulin2016fasttext}. For its initial configuration, the classifier is trained on positive samples from OpenWebMath \citep{openwebmath}, while a range of unrelated web pages function as negative examples. After training, the classifier is leveraged to extract further positive samples from the CC, which are then curated and verified by human annotators. The updated set is used to retrain the classifier, thereby enhancing its discriminative power. Subsequent evaluations demonstrate the superior quality of the large-scale dataset: our foundational model, DeepSeekMath-Base 7B, attains scores of 64.2\% on GSM8K \citep{gsm8k} and 36.2\% on the MATH competition dataset \citep{MATH}, clearly exceeding the performance of Minerva 540B \citep{minerva}. Additionally, the multilingual nature of the DeepSeekMath Corpus leads to notable gains on Chinese math benchmarks \citep{wei2023cmath,agieval}. We view our methodology in curating mathematical data as an initial contribution to the field, acknowledging ample opportunities for further advancements.

DeepSeekMath-Base is built upon DeepSeek-Coder-Base-v1.5 7B \citep{deepseek-coder}, given our findings that initializing from a code-pretrained model yields superior results over starting with a standard LLM. Additionally, we find that mathematical training not only bolsters the model's performance on tasks such as MMLU \citep{mmlu} and BBH \citep{bbh}, but also broadens its general reasoning proficiency beyond mathematics.

Following the pre-training phase, we conduct mathematical instruction tuning on DeepSeekMath-Base, utilizing datasets featuring chain-of-thought \citep{cot}, program-of-thought \citep{pot,pal}, and tool-integrated reasoning \citep{tora} approaches. The tuned model, DeepSeekMath-Instruct 7B, surpasses all other models of the same 7B parameter scale and achieves performance on par with 70B open-source models that have undergone instruction tuning.

Additionally, we present Group Relative Policy Optimization (GRPO), a modified reinforcement learning (RL) technique derived from Proximal Policy Optimization (PPO) \citep{schulman2017proximal}. Unlike traditional PPO, GRPO eliminates the need for a critic network by estimating baselines directly from group-level scores, thereby greatly lowering computational costs during training. Utilizing only a fraction of English instruction-tuning data, GRPO achieves marked performance gains over the already strong DeepSeekMath-Instruct baseline, delivering notable improvements on both in-domain tasks (GSM8K: 82.9\% $\rightarrow$ 88.2\%, MATH: 46.8\% $\rightarrow$ 51.7\%) and out-of-domain challenges (for instance, CMATH: 84.6\% $\rightarrow$ 88.8\%) during RL fine-tuning. We further propose a unified framework that encompasses diverse methods such as Rejection Sampling Fine-Tuning (RFT) \citep{yuan2023scaling}, Direct Preference Optimization (DPO) \citep{dpo}, PPO, and our GRPO approach, revealing that these can be interpreted as direct or simplified forms of RL strategies. Through comprehensive experimentation—including comparisons between online and offline training, outcome-based versus process-based supervision, and single-step versus iterative RL—we systematically examine the fundamental characteristics underlying this paradigm. Finally, we discuss the mechanisms through which RL enhances the capabilities of instruction-tuned models and outline avenues for advancing RL effectiveness within this unified framework.

\subsection{Contributions}
We present a scalable approach to math pre-training, complemented by an in-depth investigation and evaluation of reinforcement learning techniques.

\textbf{Math Pre-Training at Scale}

\begin{itemize}[topsep=0pt]
    \item Our findings demonstrate that the openly available Common Crawl dataset holds significant resources pertinent to mathematics. Utilizing a carefully crafted data selection process, we developed the DeepSeekMath Corpus, a superior dataset comprised of 120B tokens extracted from web pages specifically screened for mathematical relevance. This dataset is nearly sevenfold larger than the collection of math web pages employed by Minerva \citep{minerva}, and approximately nine times greater than the newly released OpenWebMath \citep{openwebmath}.

\item DeepSeekMath-Base 7B, our foundational pre-trained model, attains results on par with Minerva 540B \citep{minerva}, demonstrating that model size alone does not determine mathematical reasoning ability. These findings suggest that even models with fewer parameters can perform well if pre-trained with high-quality data.

\item We present results from our experiments involving math training. Pretraining models on code before engaging in math-focused training enhances their capacity to tackle mathematical problems, regardless of whether external tools are used. This provides preliminary evidence toward resolving the debated issue: \textit{does code training enhance reasoning skills?} Our evidence suggests that it does, specifically with respect to mathematical reasoning.

\item While it is standard practice to train models using arXiv papers—particularly within mathematics-focused research—this approach does not yield significant gains across any of the mathematical benchmarks considered in this work.
\end{itemize}

\textbf{Exploration and Analysis of Reinforcement Learning}

\begin{itemize}[topsep=0pt]
    \item We present Group Relative Policy Optimization (GRPO), a novel reinforcement learning approach that operates without a critic model. By estimating the baseline using group scores, GRPO dramatically lowers the computational cost of training compared to Proximal Policy Optimization (PPO).
    \item Our results show that GRPO markedly improves the capabilities of our instruction-tuned model, DeepSeekMath-Instruct, using only instruction-tuning datasets. Additionally, the method yields gains in out-of-domain generalization during reinforcement learning.
    \item We develop a comprehensive framework that contextualizes a variety of techniques—including RFT, DPO, PPO, and GRPO—under a shared perspective. In-depth empirical analyses are carried out, examining aspects such as online versus offline learning, differences between outcome and process supervision, and contrasts between single-turn and iterative reinforcement learning, to thoroughly elucidate the fundamental components within this unified framework.
    \item Leveraging this unified paradigm, we analyze the underlying mechanisms contributing to reinforcement learning’s effectiveness, and identify multiple promising avenues for advancing the reinforcement learning of large language models (LLMs).
\end{itemize}

\subsection{Summary of Evaluations and Metrics}

\begin{itemize}[topsep=0pt]
    \item \textbf{English and Chinese Mathematical Reasoning}: We systematically evaluate our models on both English and Chinese mathematical reasoning benchmarks, which span topics from elementary through college-level mathematics. The English datasets examined are GSM8K \citep{gsm8k}, MATH \citep{MATH}, SAT \citep{llemma}, OCW Courses \citep{minerva}, and MMLU-STEM \citep{mmlu}. The Chinese evaluations utilize MGSM-zh \citep{mgsm}, CMATH \citep{wei2023cmath}, Gaokao-MathCloze \citep{agieval}, and Gaokao-MathQA \citep{agieval}. Our analysis includes both the capability of models to produce standalone, tool-free textual solutions and their proficiency in employing Python code for problem-solving.

On various English evaluation benchmarks, DeepSeekMath-Base achieves results comparable to the proprietary Minerva 540B model \citep{minerva}, and clearly outperforms all open-source base models—including Mistral 7B \citep{mistral} and Llemma-34B \citep{llemma}—regardless of whether or not these models were specifically pre-trained on mathematical data, often by a wide margin. Importantly, DeepSeekMath-Base demonstrates an even greater advantage on Chinese benchmarks. This superior performance can be attributed to our approach, which, in contrast to prior studies \citep{minerva,llemma}, does not restrict math pre-training data to English only, but incorporates high-quality mathematical content from non-English sources as well. Furthermore, through the application of mathematical instruction tuning and reinforcement learning, DeepSeekMath-Instruct and DeepSeekMath-RL achieve impressive results, surpassing the 50\% accuracy threshold for the first time on the competition-level MATH dataset within the open-source landscape.

\item \textbf{Formal Mathematics}: We assess DeepSeekMath-Base on the informal-to-formal theorem proving benchmark proposed in \citep{dsp_proof} utilizing miniF2F \citep{minif2f}, with Isabelle \citep{isabelle} serving as the designated proof assistant. Our results indicate that DeepSeekMath-Base achieves notable few-shot performance in autoformalization.
\item \textbf{Natural Language Understanding, Reasoning, and Code}: To thoroughly examine the breadth of general comprehension, reasoning, and programming proficiency, we subject DeepSeekMath-Base to the Massive Multitask Language Understanding (MMLU) benchmark \citep{mmlu}, comprising 57 multiple-choice questions across a range of topics, as well as BIG-Bench Hard (BBH) \citep{bbh}, which features 23 demanding problems that generally require multi-step inferential skills. Additionally, HumanEval \citep{codex} and MBPP \citep{mbpp}, standard assessments in code generation, are used. Pre-training with mathematical data is shown to enhance the model's overall language understanding and reasoning effectiveness.

\section{Math Pre-Training}

\subsection{Data Collection and Decontamination} This section describes the methodology employed to build the DeepSeekMath Corpus utilizing data from Common Crawl. Figure \ref{fig:data_collect} illustrates a stepwise pipeline used to progressively curate an extensive mathematical dataset, beginning with a seed corpus—namely, a modest yet high-caliber math-focused dataset. Importantly, the same methodology is adaptable for constructing datasets in other fields, including programming.

Initially, we select OpenWebMath \citep{openwebmath}, a curated dataset composed of high-quality mathematical web content, as our foundational seed corpus. Leveraging this collection, we train a fastText model \citep{joulin2016fasttext} with the objective of retrieving additional mathematical web pages similar in style to those in OpenWebMath. For this, 500,000 entries are randomly drawn from the seed corpus to serve as positive samples, while 500,000 pages sourced from Common Crawl are utilized as negative counterparts. Training is carried out via a publicly available library\footnote{\url{https://fasttext.cc}}, setting the embedding dimension at 256, using a learning rate of 0.1, restricting the maximum word n-gram size to 3, ensuring a minimum occurrence threshold of 3 for words, and running for 3 epochs. In order to efficiently handle the massive scale of Common Crawl, we apply both URL-based deduplication and near-duplicate removal, ultimately distilling the collection down to 40 billion unique HTML pages. The trained fastText model is then used to surface mathematical web pages from this deduplicated dataset. To ensure content quality, we sort the retrieved pages based on the fastText model’s score, retaining only the highest-ranked items. Data retention levels are subsequently validated through pre-training experiments across data slices consisting of the top 40B, 80B, 120B, and 160B tokens, with the first phase prioritizing retention of the top 40B tokens.

Following the initial round of data acquisition, a substantial portion of mathematical web pages remains unretrieved, primarily due to the limited diversity in the positive training examples for the fastText model. To address this issue, we enhance the seed corpus by identifying further web sources focused on mathematics, aiming to refine the fastText model’s effectiveness. Our methodology begins with partitioning the complete Common Crawl dataset into non-overlapping domains, where a domain refers to all web pages that possess the same base URL. For every domain, we compute the proportion of pages retrieved in the first collection phase. Domains in which over 10\% of web pages are gathered are deemed to be associated with mathematics (for example, \url{mathoverflow.net}). 

We then proceed to manually curate URLs that are linked with mathematical content inside these domains (such as \url{mathoverflow.net/questions}). Those web pages connected to the identified URLs but not yet acquired are incorporated into the seed corpus. By expanding the set of positive instances in this way, we facilitate the training of a more effective fastText model that can capture a greater volume of mathematical data in future iterations. Across four cycles of data gathering, our collection grows to 35.5M mathematical web pages, amounting to a total of 120B tokens. In the fourth iteration, we observe that approximately 98\% of the material had already been accumulated in the preceding iteration, prompting us to discontinue further data collection.

In order to prevent contamination of benchmarks, we adopt the methodology described by \cite{deepseek-coder}, eliminating web content that features questions or answers sourced from English math benchmarks such as GSM8K \citep{gsm8k} and MATH \citep{MATH}, as well as Chinese datasets like CMATH \citep{wei2023cmath} and AGIEval \citep{agieval}. Specifically, our filtering approach excludes from the math training dataset any segment of text containing a 10-gram string that exactly matches any substring present in these benchmark resources. For benchmark passages that are shorter than 10 grams but consist of at least 3 grams, we apply strict exact matching to identify and remove affected web pages from the corpus.

\subsection{Validating the Quality of the DeepSeekMath~Corpus}

We run pre-training experiments to investigate how the DeepSeekMath~Corpus is compared with the recently released math-training corpora:

\begin{itemize}[topsep=0pt]
    \item \textbf{MathPile} \citep{mathpile}: a composite dataset containing 8.9B tokens, compiled from diverse sources such as textbooks, Wikipedia, ProofWiki, CommonCrawl, StackExchange, and arXiv; notably, more than 85\% of the material is derived from arXiv;
    \item \textbf{OpenWebMath} \citep{openwebmath}: a subset of CommonCrawl curated specifically for mathematics-related content, comprising 13.6B tokens;
    \item \textbf{Proof-Pile-2} \citep{llemma}: a dataset constructed from OpenWebMath, AlgebraicStack (representing 10.3B tokens of mathematics code), and 28.0B tokens of arXiv articles. In experiments using Proof-Pile-2, we adopt the arXiv:Web:Code proportional split of 2:4:1 as established in \cite{llemma}.
\end{itemize}

\subsubsection{Training Setting}
\label{sec:quality-policy}
We fine-tune a general-purpose pre-trained language model comprising 1.3B parameters, structurally equivalent to the DeepSeek LLMs \citep{deepseek-llm}, referred to as DeepSeek-LLM 1.3B in this work. For each mathematical dataset, we train a separate model over 150B tokens. All training processes utilize the streamlined HAI-LLM framework \citep{haillm}, prioritizing efficiency and scalability. Consistent with the training regimen described for DeepSeek LLMs, we implement the AdamW optimizer \citep{adamW} with hyperparameters $\beta_1=0.9$, $\beta_2=0.95$, and a $\mathrm{weight\_decay}$ of 0.1. The learning rate is governed by a multi-stage scheduler: it linearly ramps up to its maximum after an initial 2,000-step warmup, drops to 31.6\% of this peak at the 80\% completion point of training, and further reduces to 10.0\% of the peak by the end of 90\% of total steps. The highest learning rate employed is 5.3e-4. Training is conducted with a batch size of 4 million tokens, each sequence accommodating up to 4,000 tokens in context.

\subsubsection{Evaluation Results}

\textbf{The DeepSeekMath~Corpus is of high quality, covers multilingual mathematical content, and is the largest in size.}

\begin{itemize}[topsep=0pt]
    \item \textbf{High-quality}:  
    We assess the performance of models on eight mathematical evaluation datasets utilizing few-shot chain-of-thought prompting as described in \cite{cot}. According to Table \ref{tab:corpora_comparison}, models trained with the DeepSeekMath~Corpus consistently outperform those trained on alternative datasets. Additionally, Figure \ref{fig:corpora_comparisons} reveals that at 50B tokens—equivalent to one complete pass over Proof-Pile-2—the model exposed to DeepSeekMath~Corpus achieves noticeably higher results compared to Proof-Pile-2, suggesting that the DeepSeekMath~Corpus offers superior average data quality.
    \item \textbf{Multilingual}:  
    The DeepSeekMath~Corpus comprises data sources in several languages, with English and Chinese constituting the majority of the material. Training on this corpus, as illustrated by Table \ref{tab:corpora_comparison}, substantially improves the ability of models to reason mathematically in both English and Chinese. By contrast, most currently available mathematical datasets are primarily focused on English, which results in limited progress and, at times, even degrades performance in mathematical reasoning tasks carried out in Chinese.
    \item \textbf{Large-scale}:  
    Compared to other mathematical datasets, the DeepSeekMath~Corpus is significantly more expansive. Figure \ref{fig:corpora_comparisons} demonstrates that the DeepSeek-LLM 1.3B, when trained on DeepSeekMath~Corpus, achieves both a sharper trajectory of improvement and more durable gains in performance. Conversely, existing baseline corpora, which are considerably smaller and frequently cycled through multiple times during training, tend to hit their peak performance early and plateau thereafter.
\end{itemize}

\subsection{Training and Evaluating DeepSeekMath-Base 7B}

This section presents DeepSeekMath-Base 7B, a foundational model designed for advanced mathematical reasoning. The model builds upon DeepSeek-Coder-Base-v1.5 7B \citep{deepseek-coder} as its starting point, and undergoes training on a corpus totaling 500B tokens. The dataset composition consists of 56\% DeepSeekMath Corpus, 4\% AlgebraicStack, 10\% arXiv, 20\% source code from Github, with the final 10\% comprised of natural language samples in both English and Chinese derived from Common Crawl. The majority of the training configuration aligns with that described in Section \ref{sec:quality-policy}, with two exceptions: the peak learning rate is established at 4.2e-4, and the batch size is increased to 10M tokens.

We thoroughly evaluate DeepSeekMath-Base 7B’s mathematical proficiency, analyzing its performance in generating independent mathematical solutions without external assistance, utilizing tools for mathematical problem-solving, and executing formal theorem proofs. In addition to these mathematical tasks, we further examine the broader attributes of the base model, such as its competence in natural language comprehension, logical reasoning, and programming capabilities.

\paragraph{Mathematical Problem Solving with Step-by-Step Reasoning}

We assess the mathematical problem-solving capabilities of DeepSeekMath-Base employing few-shot chain-of-thought prompting \citep{cot} over eight English and Chinese benchmarks. These benchmarks include both quantitative reasoning tasks (such as GSM8K \citep{gsm8k}, MATH \citep{MATH}, and CMATH \citep{wei2023cmath}) and multiple-choice formats (including MMLU-STEM \citep{mmlu} and Gaokao-MathQA \citep{agieval}), spanning a broad spectrum of mathematical domains and ranging in difficulty from elementary through college-level.

According to Table \ref{tab:base_math_cot}, DeepSeekMath-Base 7B consistently achieves the highest scores across all eight benchmarks when compared to other open-source base models, such as the commonly adopted Mistral 7B \citep{mistral} and the newly introduced Llemma 34B \citep{llemma}, which benefited from math-focused training on Proof-Pile-2 \citep{llemma}. Remarkably, on the competition-grade MATH dataset, DeepSeekMath-Base outperforms all available open-source base models by a margin exceeding 10\% in absolute terms, and even surpasses Minerva 540B \citep{minerva}, a proprietary model with 77 times more parameters, built atop PaLM \citep{palm} and further refined on mathematical corpora.

\paragraph{Mathematical Problem Solving with Tool Use} We assess the capability of models to solve mathematical problems on the GSM8K and MATH benchmarks by employing few-shot program-of-thought prompting \citep{pot,pal}. In this framework, models are asked to generate Python code to solve each problem, with access to libraries like \textit{math} and \textit{sympy} to handle complex computations. The correctness of the program's output is used to determine the answer. As indicated in Table \ref{tab:base_math_pot_proof}, DeepSeekMath-Base 7B achieves superior results compared to the previous leading model, Llemma 34B.

\paragraph{Formal Mathematics}

Automating formal proofs plays a significant role in verifying the correctness and dependability of mathematical arguments while also improving productivity, and has garnered increasing interest recently. In this study, we assess the capabilities of DeepSeekMath-Base 7B for the informal-to-formal proving task as described in \citep{dsp_proof}. This task requires the generation of a formal proof from an informal statement, its formal equivalent, and an informal proof. The evaluation uses the miniF2F \citep{minif2f} benchmark, which focuses on formalizing mathematics at the Olympiad level. For each problem, we employ few-shot prompting to produce a formal Isabelle proof. Adhering to the approach in \cite{dsp_proof}, we have the models create proof sketches, after which the automated proving tool Sledgehammer \citep{sledgehammer} is utilized to complete the proofs by providing missing steps. As indicated in Table \ref{tab:base_math_pot_proof}, DeepSeekMath-Base 7B achieves notably strong results in the autoformalization of proofs.

\paragraph{Natural Language Understanding, Reasoning, and Code}

We assess the natural language understanding ability of the models using MMLU \citep{mmlu}, reasoning skills with BBH \citep{bbh}, and coding performance via HumanEval \citep{codex} and MBPP \citep{mbpp}. According to Table \ref{tab:understanding_reasoning_code}, DeepSeekMath-Base 7B demonstrates marked improvements over the prior DeepSeek-Coder-Base-v1.5 \citep{deepseek-coder} on both MMLU and BBH, highlighting the beneficial influence of mathematics-oriented training on both language comprehension and reasoning tasks. Moreover, by incorporating code tokens during continued training, DeepSeekMath-Base 7B successfully preserves the coding benchmark results achieved by DeepSeek-Coder-Base-v1.5. In summary, DeepSeekMath-Base 7B substantially surpasses the general-purpose Mistral 7B \citep{mistral} across the evaluated reasoning and coding tasks.

\section{Supervised Fine-Tuning}

\subsection{SFT Data Curation}

A diverse instruction-tuning dataset for mathematics is assembled, incorporating problems in both English and Chinese from a range of mathematical disciplines and encompassing multiple levels of difficulty. Each problem is annotated with corresponding solutions presented in chain-of-thought (CoT) \citep{cot}, program-of-thought (PoT) \citep{pot,pal}, and tool-integrated reasoning \citep{tora} formats. The resulting training corpus consists of 776K examples.

\begin{itemize}[topsep=0pt]
    \item \textbf{English mathematical datasets}: Our annotation encompasses GSM8K and MATH with solutions that incorporate tool usage. We further include a selected portion of MathInstruct \citep{MathInstruct} and utilize the training split of Lila-OOD \citep{lila}, both featuring problems addressed via CoT or PoT methods. This English dataset suite spans various branches of mathematics, including algebra, probability, number theory, calculus, and geometry.
    \item \textbf{Chinese mathematical datasets}: We assemble a collection of Chinese K-12 mathematics problems, covering 76 sub-domains such as linear equations. Each problem is annotated with reasoning presented through both CoT and tool-integrated approaches.
\end{itemize}

\subsection{Training and Evaluating DeepSeekMath-Instruct 7B}

This section presents DeepSeekMath-Instruct 7B, which is developed by fine-tuning DeepSeekMath-Base with mathematical instructions. For data preparation, training samples are combined at random until the sequence length cap of 4K tokens is met. The model is optimized over 500 training steps, utilizing a batch size of 256 and maintaining a fixed learning rate of 5e-5.

We evaluate models' mathematical performance both without and with tool use, on 4 quantitative reasoning benchmarks in English and Chinese. We benchmark our model against the leading models of the time:

\begin{itemize}[topsep=0pt]
    \item \textbf{Closed-source models} encompass the following: (1) the GPT series, with GPT-4 \citep{gpt4} and its Code Interpreter~\footnote{\url{https://openai.com/blog/chatgpt-plugins\#\#code-interpreter}} being among the most advanced, (2) Gemini Ultra and Pro \citep{gemini}, (3) Inflection-2 \citep{inflection-2}, (4) Grok-1~\footnote{\url{https://x.ai/model-card}},
    along with newly introduced Chinese models such as (5) Baichuan-3~\footnote{\url{https://www.baichuan-ai.com}},
    and (6) the most recent version of GLM-4~\footnote{\url{https://open.bigmodel.cn/dev/api\#glm-4}}.

    from the GLM family \citep{glm}.

These models are intended for broad application, and the majority have been refined through various alignment methods.

\item \textbf{Open-source models} feature a range of general-purpose systems, such as (1) DeepSeek-LLM-Chat 67B \citep{deepseek-llm}, (2) Qwen 72B \citep{qwen}, (3) SeaLLM-v2 7B \citep{seallm}, and (4) ChatGLM3 6B \citep{chatglm3}. In addition, there are models with specialized improvements in mathematical reasoning, including: (5) InternLM2-Math 20B~\footnote{\url{https://github.com/InternLM/InternLM-Math}}, which extends InternLM2 by incorporating targeted math training and subsequent instruction tuning; (6) Math-Shepherd-Mistral 7B, in which Mistral 7B \citep{mistral} is optimized with PPO training \citep{schulman2017proximal} via a process-supervised reward mechanism; (7) the WizardMath collection \citep{wizardmath}, which augments mathematical reasoning abilities in Mistral 7B and Llama-2 70B \citep{llama2} through evolve-instruct—a form of instruction tuning that utilizes AI-generated instructions—and PPO training, using datasets primarily from GSM8K and MATH; (8) MetaMath 70B \citep{metamath}, a fine-tuned Llama-2 70B trained on a synthesized version of GSM8K and MATH; (9) ToRA 34B \cite{tora}, developed by further fine-tuning CodeLlama 34B for mathematical reasoning with integrated tool use; and (10) MAmmoTH 70B \citep{MathInstruct}, a Llama-2 70B model instruction-tuned on MathInstruct.

Referencing Table \ref{tab:sft_rl_math}, when evaluating without permitting tool use, DeepSeekMath-Instruct 7B exhibits robust step-by-step reasoning abilities. Specifically, on the MATH dataset designed for competition-level tasks, our model exceeds the accuracy of every open-source competitor as well as most proprietary systems—including Inflection-2 and Gemini Pro—by at least a 9\% absolute margin. This holds even when compared to significantly larger models (such as Qwen 72B) or those benefiting from targeted reinforcement learning focused on mathematics (like WizardMath-v1.1 7B). Although DeepSeekMath-Instruct achieves parity with Chinese proprietary models like GLM-4 and Baichuan-3 on the MATH benchmark, it does not yet match the performance levels of GPT-4 or Gemini Ultra.

When models are permitted to combine natural language reasoning with programmatic tool use during evaluation, DeepSeekMath-Instruct 7B achieves nearly 60\% accuracy on the MATH benchmark, outperforming all previously released open-source models. Across additional benchmarks, our model demonstrates performance comparable to DeepSeek-LLM-Chat 67B—the earlier state-of-the-art—which is ten times its size.
\section{Reinforcement Learning}

\subsection{Group Relative Policy Optimization} After the Supervised Fine-Tuning (SFT) phase, reinforcement learning (RL) has demonstrated its capability to enhance the mathematical reasoning performance of LLMs \citep{wang2023math,wizardmath}. Here, we present Group Relative Policy Optimization (GRPO), an RL method we have developed that is both efficient and effective.

\subsubsection{From PPO to GRPO} Proximal Policy Optimization (PPO) \citep{schulman2017proximal} is an actor-critic RL algorithm that is widely used in the RL fine-tuning stage of LLMs \citep{ouyang2022training}. In particular, it optimizes LLMs by maximizing the following surrogate objective:

\begin{equation}
\footnotesize
    \mathcal{J}_{PPO}(\theta) = \mathbb{E}{[q \sim P(Q), o \sim \pi_{\theta_{old}}(O|q)]} \frac{1}{|o|} \sum_{t=1}^{|o|} \min \left[ \frac{\pi_\theta(o_{t} | q, o_{<t})}{\pi_{\theta_{old}}(o_{t} | q, o_{<t})} A_{t}, \text{clip} \left( \frac{\pi_\theta(o_{t} | q, o_{<t})}{\pi_{\theta_{old}}(o_{t} | q, o_{<t})}, 1 - \varepsilon, 1 + \varepsilon \right)  A_{t} \right] ,
\end{equation}

Here, $\pi_{\theta}$ denotes the updated policy model, while $\pi_{\theta_{old}}$ corresponds to the previous policy version. The variables $q$ and $o$ refer to the sampled questions from the dataset and responses generated by $\pi_{\theta_{old}}$, respectively. The parameter $\varepsilon$ represents a clipping hyper-parameter used in PPO to enhance training stability. The term $A_t$ signifies the advantage, derived using Generalized Advantage Estimation (GAE) \citep{gae} and calculated from subsequent rewards $\{r_{\ge t}\}$ as well as an estimated value function $V_{\psi}$. Consequently, PPO requires joint training of the policy and value function. To counteract the issue of excessive reward model optimization, it is conventional to introduce a per-token KL-divergence penalty computed from a reference model into the reward at each token step, as proposed by \citep{ouyang2022training}; that is,

\begin{equation}
    r_{t} = r_\varphi(q, o_{\le t}) - \beta \log\frac{\pi_{\theta}(o_{t}|q, o_{<t})}{\pi_{ref}(o_{t}|q, o_{<t})},
\label{eq:PPO-reward}
\end{equation}

Here, $r_\varphi$ denotes the reward model, while $\pi_{ref}$ refers to the reference model—commonly chosen as the initial SFT model—and $\beta$ represents the weight assigned to the KL penalty.

As the value function employed in PPO is typically another model of comparable size as the policy model, it brings a substantial memory and computational burden. Additionally, during RL training, the value function is treated as a baseline in the calculation of the advantage for variance reduction. While in the LLM context, usually only the last token is assigned a reward score by the reward model, which may complicate the training of a value function that is accurate at each token. To address this, as shown in Figure \ref{fig:grpo}, we propose Group Relative Policy Optimization (GRPO), which obviates the need for additional value function approximation as in PPO, and instead uses the average reward of multiple sampled outputs, produced in response to the same question, as the baseline. More specifically, for each question $q$, GRPO samples a group of outputs $\{o_1, o_2, \cdots, o_G\}$  from the old policy  $\pi_{\theta_{old}}$  and then optimizes the policy model by maximizing the following objective:

\begin{equation}
\footnotesize
\begin{split}
    \mathcal{J}_{GRPO}(\theta) &= \mathbb{E}_{q \sim P(Q), \{o_i\}_{i=1}^G \sim \pi_{\theta_{old}}(O|q)}  \\
    & \frac{1}{G}\sum_{i=1}^G\frac{1}{|o_i|} \sum_{t=1}^{|o_i|} \Bigg\{ \min \Bigg[ \frac{\pi_\theta(o_{i,t} | q, o_{i,<t})}{\pi_{\theta_{old}}(o_{i,t} | q, o_{i,<t})} \hat{A}_{i,t}, \; \text{clip} \Bigg( \frac{\pi_\theta(o_{i,t} | q, o_{i,<t})}{\pi_{\theta_{old}}(o_{i,t} | q, o_{i,<t})}, \; 1 - \varepsilon, \; 1 + \varepsilon \Bigg) \hat{A}_{i,t} \Bigg] - \beta \mathbb{D}_{KL}\left[\pi_{\theta} || \pi_{ref}\right] \Bigg\},
\end{split}
\label{eq:GRPO-obj}
\end{equation}

where $\varepsilon$ and $\beta$ are hyper-parameters, and $\hat{A}_{i,t}$  is the advantage calculated based on relative rewards of the outputs inside each group only, which will be detailed in the following subsections. The group relative way that GRPO leverages to calculate the advantages, aligns well with the comparative nature of rewards models,  as reward models are typically trained on datasets of comparisons between outputs on the same question. Also note that, instead of adding KL penalty in the reward, GRPO regularizes by directly adding the KL divergence between the trained policy and the reference policy to the loss, avoiding complicating the calculation of  $\hat{A}_{i,t}$. And different from the KL penalty term used in (\ref{eq:PPO-reward}), we estimate the KL divergence with the following unbiased estimator \citep{kl_approx}:

\begin{equation}
\small
    \mathbb{D}_{KL}\left[\pi_{\theta} || \pi_{ref}\right] = \frac{\pi_{ref}(o_{i,t}|q,o_{i,<t})}{\pi_{\theta}(o_{i,t}|q,o_{i,<t})} - \log \frac{\pi_{ref}(o_{i,t}|q,o_{i,<t})}{\pi_{\theta}(o_{i,t}|q,o_{i,<t})} - 1,
\end{equation}

which is guaranteed to be positive.

\begin{algorithm}[t]
  \small
  \caption{Iterative Group Relative Policy Optimization}
  \textbf{Input} initial policy model $\pi_{\theta_{\text{init}}}$; reward models $r_\varphi$; task prompts $\mathcal{D}$; 
  hyperparameters $\varepsilon$, $\beta$, $\mu$
  \begin{algorithmic}[1]
    \State Initialize policy model $\pi_\theta$ as $\pi_{\theta_{\text{init}}}$
    \For{iteration = 1, \dots, I}
       \State Set reference model $\pi_{ref}$ to current $\pi_{\theta}$
      \For{step = 1, \dots, M}
      \State Draw a minibatch $\mathcal{D}_b$ from $\mathcal{D}$
      \State Assign $\pi_{\theta_{old}}$ to current policy $\pi_{\theta}$
      \State For each question $q$ in $\mathcal{D}_b$, generate $G$ outputs $\{o_i\}_{i=1}^G$ by sampling from $\pi_{\theta_{old}} (\cdot \mid q)$
      \State Use the reward model $r_{\varphi}$ to determine the set of rewards $\{r_i\}_{i=1}^G$ corresponding to each output $o_i$
      \State Estimate group relative advantages $\hat{A}_{i,t}$ for the $t$-th token in every $o_i$
      \For{GRPO iteration = 1, \dots, $\mu$}
        \State Adjust $\pi_{\theta}$ by optimizing the GRPO objective (Equation \ref{eq:GC-GRPO})
      \EndFor
    \EndFor \State Update the reward model $r_\varphi$ via continual training with a replay strategy.
    \EndFor 
  \end{algorithmic}
  \textbf{Output} $\pi_\theta$
  \label{alg:iter-grpo}
\end{algorithm}

\subsubsection{Outcome Supervision RL with GRPO} Consider a given question $q$, for which a batch of outputs $\{o_1, o_2, \cdots, o_G\}$ is generated via the previous policy $\pi_{\theta_{old}}$. Each of these outputs is then evaluated by a reward model, resulting in a set of rewards $\mathbf{r}=\{r_1, r_2, \cdots, r_G\}$. These reward values undergo normalization—specifically, subtracting the mean of the group and dividing by their standard deviation. Within the outcome supervision framework, the normalized reward associated with each output $o_i$ is assigned as the advantage $\hat{A}_{i, t}$ for every token in that output, namely $\hat{A}_{i, t} = \widetilde{r}_i = \frac{r_i- {\rm mean}(\mathbf{r})}{{\rm std}(\mathbf{r})}$. The policy is then trained to maximize the objective stated in equation (\ref{eq:GRPO-obj}).

\subsubsection{Process Supervision RL with GRPO}

In outcome supervision, the model only receives feedback upon the completion of its generated sequence, which often fails to deliver adequate and timely guidance, particularly for intricate mathematical problems. Inspired by \cite{wang2023math}, we further investigate process supervision, wherein rewards are assigned after each individual reasoning step rather than only at the final output. Specifically, for a given question $q$ and $G$ sampled completions $\{o_1, o_2, \cdots, o_G\}$, a process reward model evaluates each reasoning step, assigning a series of step-level rewards: $\mathbf{R} = \{ \{r_1^{index(1)},\cdots,r_1^{index(K_1)}\}, \cdots,  \{r_G^{index(1)},\cdots,r_G^{index(K_G)}\} \}$. Here, $index(j)$ indicates the index of the end token for the $j$-th step, and $K_i$ denotes the number of steps in the $i$-th sampled output. These raw rewards are then standardized by subtracting their mean and dividing by their standard deviation, so that $\widetilde{r}_i^{index(j)} = \frac{r_i^{index(j)} - {\rm mean(\mathbf{R})}}{{\rm std(\mathbf{R})}}$.

With this setup, for each token, the advantage is computed as the sum of all normalized rewards of the steps that follow (including itself), i.e., $\hat{A}_{i, t} = \sum_{index(j) \ge t} \widetilde{r}_i^{index(j)}$. This advantage signal is then employed to optimize the policy through maximization of the objective presented in equation (\ref{eq:GRPO-obj}).

\subsubsection{Iterative RL with GRPO} During reinforcement learning, the existing reward model can become inadequate for effectively guiding the updated policy model. To address this, we investigate an iterative RL approach using GRPO. As outlined in Algorithm \ref{alg:iter-grpo}, iterative GRPO involves producing fresh datasets for the reward model, derived from samples generated by the policy model. The reward model is incrementally updated through continual training that employs a replay mechanism, which preserves 10\% of earlier data. Subsequently, the policy model serves as the new reference model, and training of the policy model proceeds with the updated reward model.

\subsection{Training and Evaluating DeepSeekMath-RL}

Our reinforcement learning experiments utilize DeepSeekMath-Instruct 7B as the policy model. For RL training, we curate data comprising approximately 144K chain-of-thought-style questions derived from the GSM8K and MATH subsets within the SFT corpus. By intentionally omitting other SFT items during the RL stage, we aim to assess how reinforcement learning influences benchmarks for which there is an absence of dedicated RL training data. The reward model's training dataset is constructed in accordance with the methodology described in \citep{wang2023math}. We initiate reward model training with DeepSeekMath-Base 7B, employing a learning rate of $2 \times 10^{-5}$. For GRPO, the policy model utilizes a learning rate of $1 \times 10^{-6}$, while the KL penalty is set to 0.04. A total of $64$ responses are sampled for each question. We set both the maximum input length and the training batch size to 1024. After every exploration cycle, the policy model undergoes only a single update. The evaluation protocol for DeepSeekMath-RL 7B mirrors that of DeepSeekMath-Instruct 7B. Here, tasks involving GSM8K and MATH datasets with chain-of-thought prompting are designated as in-domain, whereas all other benchmarks fall under the category of out-of-domain evaluation.

Table \ref{tab:sft_rl_math} presents comparative results for both open- and closed-source models, evaluating their abilities in chain-of-thought and tool-integrated reasoning across English and Chinese datasets. Our analysis reveals that: 1) DeepSeekMath-RL 7B achieves 88.2\% accuracy on GSM8K and 51.7\% on MATH benchmarks when employing chain-of-thought reasoning, outperforming every open-source model within the 7B to 70B parameter category, as well as most closed-source systems. 2) Notably, DeepSeekMath-RL 7B is exclusively fine-tuned on chain-of-thought-style instruction data from GSM8K and MATH, beginning with DeepSeekMath-Instruct 7B as the base model. Even though its training corpus is relatively limited, DeepSeekMath-RL 7B demonstrates superior performance to DeepSeekMath-Instruct 7B on every metric evaluated, illustrating the substantial benefits conferred by reinforcement learning.

\section{Discussion}
This section presents an overview of our results obtained from both pre-training and RL experiments.

\subsection{Lessons Learnt in Pre-Training}

To begin, we discuss our pre-training experience. Except where otherwise noted, our training procedures follow the configuration described in Section \ref{sec:quality-policy}. Additionally, references to the DeepSeekMath Corpus in this section pertain to an 89B-token dataset derived from the second round of data gathering.

\subsubsection{Code Training Benefits Mathematical Reasoning}

An often-cited but unproven theory posits that practicing coding enhances reasoning skills. In this work, we seek to address this claim, specifically as it pertains to mathematics: training on code leads to enhanced mathematical reasoning in models, regardless of whether external tools are employed.

To study how code training affects mathematical reasoning, we experimented with the following two-stage training and one-stage training settings:

\textbf{Two-Stage Training}

\begin{itemize}[topsep=0pt]
    \item \textbf{Code Training for 400B Tokens $\rightarrow$ Math Training for 150B Tokens}: Our approach involves initially pre-training DeepSeek-LLM 1.3B on 400B code tokens, subsequently continuing training with 150B math tokens;
    \item \textbf{General Training for 400B Tokens $\rightarrow$ Math Training for 150B Tokens}: For comparison, we conduct a parallel experiment where the first 400B tokens used for pre-training are drawn from a broad general-purpose corpus assembled by DeepSeek-AI, aiming to assess whether code token pre-training offers specific benefits for enhancing mathematical reasoning when compared with general token pre-training.
\end{itemize}

\textbf{One-Stage Training}

\begin{itemize}[topsep=0pt]
    \item \textbf{Math Training with 150B Tokens}: We utilize a corpus of 150B mathematical tokens to train DeepSeek-LLM 1.3B;
    \item \textbf{Joint Training on 400B Code Tokens and 150B Math Tokens}: Sequential math training after code pre-training leads to reduced coding ability. We therefore explore if integrating code and math tokens simultaneously during single-stage training can enhance mathematical reasoning skills while also mitigating catastrophic forgetting.
\end{itemize}

\paragraph{Results}
The downstream outcomes observed across various training configurations are presented in Table \ref{tab:code-math} and Table \ref{tab:code-math-general-eval}.

Training with code demonstrably enhances program-aided mathematical reasoning, regardless of whether a two-stage or one-stage training protocol is utilized. Table \ref{tab:code-math} indicates that, in the two-stage approach, code-focused training alone leads to substantial gains in the model’s capacity to solve Python-based GSM8K and MATH questions. Subsequent math training in the second phase produces additional performance gains. Notably, when code and math tokens are combined during a single-stage training process, this integration effectively prevents the catastrophic forgetting commonly observed in the two-stage framework. Furthermore, such joint training not only reinforces coding competencies (as seen in Table \ref{tab:code-math-general-eval}), but also enhances program-assisted mathematical problem solving (as demonstrated in Table \ref{tab:code-math}).

Training with code also enhances mathematical reasoning skills even when tools are not employed. In the two-stage training paradigm, beginning with code training yields noticeable gains, which in turn facilitates more effective mathematics training in the following phase, culminating in optimal overall results. Conversely, integrating code and mathematical tokens within a single-stage training framework undermines performance in mathematical reasoning tasks when tools are absent. This may be attributed to the relatively small model size of DeepSeek-LLM 1.3B, which appears insufficient for concurrently learning both code and mathematical domains.

\subsubsection{ArXiv Papers Seem Ineffective in Improving Mathematical Reasoning}

ArXiv papers are commonly included as a component of math pre-training data \citep{minerva,gpt-f,llemma,mathpile}. However, detailed analysis regarding their impact on mathematical reasoning has not been extensively conducted. Perhaps counter-intuitively, according to our experiments, arXiv papers seem ineffective in improving mathematical reasoning. We experiment with models of different sizes, including DeepSeek-LLM 1.3B and DeepSeek-Coder-Base-v1.5 7B \citep{deepseek-coder}, using arXiv corpora that underwent varied processing pipelines:

\begin{itemize}[topsep=0pt]
    \item \textbf{MathPile} \citep{mathpile}: an 8.9B-token dataset assembled through a combination of heuristic data cleaning and filtering procedures, consisting of over 85\% scientific documents from arXiv;
    \item \textbf{ArXiv-RedPajama} \citep{redpajama}: a collection comprising all arXiv LaTeX source files, where non-essential elements such as preambles, comments, macros, and bibliographies have been stripped, amounting to 28.0B tokens in total.
\end{itemize}

In our study, we independently train DeepSeek-LLM 1.3B on 150B tokens and DeepSeek-Coder-Base-v1.5 7B on 40B tokens using separate arXiv corpora. The results indicate that exposure to arXiv papers does not enhance the models' mathematical reasoning abilities. Training exclusively on arXiv content yields either negligible gains or declines in performance on a range of mathematical evaluation datasets of varying difficulty considered in this work. These evaluation sets encompass quantitative reasoning tasks such as GSM8K and MATH (Table \ref{tab:arxiv-cot}), multiple-choice assessments like MMLU-STEM (Table \ref{tab:arxiv-cot}), as well as formal mathematics datasets like miniF2F (Table \ref{tab:arxiv_atp}).

However, this conclusion has its limitations and should be taken with a grain of salt. We have not yet studied:

\begin{itemize}[topsep=0pt]
    \item How arXiv tokens influence particular mathematical tasks outside the scope of this study, for instance, the informalization of theorems, where formal statements or proofs are translated into informal language;
    \item The interaction between arXiv tokens and different categories of data;
    \item If scaling up the model would amplify the advantages offered by arXiv papers.
\end{itemize}

Thus, further exploration is required, which we leave for future studies.

\subsection{Insights of Reinforcement Learning}
\subsubsection{Towards to a Unified Paradigm} This section introduces a cohesive framework for examining various training approaches, including SFT, RFT, DPO, PPO, and GRPO. We also present experimental investigations that probe the contributing components within this unified framework. In a general sense, the parameter gradient $\theta$ under different training strategies may be expressed as follows:

\begin{equation}
    \nabla_{\theta}\mathcal{J}_{\textcolor{red}{\mathcal{A}}}(\theta) = \mathbb{E}[\underbrace{(q,o) \sim \textcolor{red}{\mathcal{D}}}_{Data \ Source}]\left( \frac{1}{|o|} \sum_{t=1}^{|o|}  \underbrace{GC_{{\mathcal{A}}}(q, o, t, \textcolor{red}{\pi_{{rf}}})}_{Gradient \ Coefficient}  \nabla_{\theta}\log \pi_{\theta}(o_t | q, o_{<t})\right).
\label{eq:objective}
\end{equation}

There exist three key components: 1) \textit{Data Source $\mathcal{D}$}, which determines the training data; 2) \textit{Reward Function $\pi_{{rf}}$}, which is the source of the training reward signal; 3) \textit{Algorithm $\mathcal{A}$}: which processes the training data and the reward signal to the gradient coefficient $GC$ that determines the magnitude of the penalty or reinforcement for the data. We analyze several representative methods based on such a unified paradigm:

\begin{itemize}
    \item  \textbf{Supervised Fine-tuning (SFT)}: SFT adapts a pretrained model by further training it on a curated dataset selected by humans.
    \item \textbf{Rejection Sampling Fine-tuning (RFT)}: RFT advances the SFT model by continuing its training using outputs generated by the SFT model in response to SFT questions, filtering these outputs to retain only those with correct answers.
    \item \textbf{Direct Preference Optimization (DPO)}: DPO further tunes the SFT model by training it with enhanced outputs from the SFT model, leveraging a pairwise DPO loss function for optimization.
    \item \textbf{Online Rejection Sampling Fine-tuning (Online RFT)}: In contrast to standard RFT, Online RFT starts with the SFT model as the initial policy and incrementally improves it through fine-tuning on augmented outputs dynamically sampled from the current policy model during training.
    \item \textbf{PPO/GRPO}: PPO/GRPO begins by initializing the policy model with the SFT model, subsequently strengthening it by sampling outputs in real time from the evolving policy model.
\end{itemize}

Table \ref{tab:data_policy} provides an overview of the elements that constitute these methods. A comprehensive derivation process can be found in Appendix \ref{app:analysis-rl}.

\paragraph{Observation about Data Source} We categorize data sources into two primary types: online sampling and offline sampling. In online sampling, training data are generated through the active exploration of the current policy model during training. In contrast, offline sampling involves collecting training data from the outputs of a pre-trained SFT model. RFT and DPO employ an offline approach, whereas Online RFT and GRPO utilize online sampling methods.

Figure \ref{fig:rl-analysis} illustrates that Online RFT consistently achieves higher performance than RFT across both benchmarks examined. While the performance of Online RFT and RFT remains similar during the initial phases of training—likely due to the actor's outputs closely matching those of the SFT model and only negligible discrepancies emerging in the sampled data—Online RFT begins to show a clear lead as training progresses. This advantage is particularly notable in the later stages, where more pronounced differences appear in the data produced by the actor, thereby underscoring the efficacy of real-time data sampling and the overall benefit of online training strategies.

\paragraph{Observation about Gradient Coefficient} Within the algorithm, the reward signal is incorporated into the calculation of the gradient coefficient, guiding the update of model parameters. In our experimental setup, we categorize the reward function into two types: `Rule' and `Model'. The `Rule' mechanism evaluates response quality strictly by answer correctness, whereas the `Model' component involves a reward model that assigns scores to responses, having been trained on data labeled by the rule-based assessment. Equations \ref{eq:GC-RFT} and \ref{eq:GC-GRPO} underscore a critical distinction between GRPO and Online RFT: GRPO specifically modulates its gradient coefficient in response to the reward value assigned by the reward model, thus enabling responses to be reinforced or penalized in proportion to the magnitude of their rewards. Conversely, Online RFT does not employ this adaptive adjustment; it omits the penalization of incorrect responses and applies equal reinforcement to all correct answers, regardless of their quality.

As illustrated in Figure \ref{fig:rl-analysis}, GRPO outperforms online RFT, underscoring the effectiveness of modifying the positive and negative gradient coefficients. Moreover, GRPO+PS achieves better results than GRPO+OS, suggesting that employing step-aware, fine-grained gradient coefficients offers distinct advantages. Additionally, our investigation of iterative RL involved conducting two successive iterations. Figure \ref{fig:iter-rl} reveals that iterative RL notably boosts performance, with the most pronounced improvement observed during the initial iteration.

\subsubsection{Why RL Works?} This study applies reinforcement learning using a selective portion of instruction tuning datasets, resulting in a marked improvement over models trained solely via instruction tuning. To investigate the underlying reasons for this effectiveness, we compare the Pass@K and Maj@K accuracies of both the Instruct and RL-enhanced models across two standard benchmarks. As depicted in Figure \ref{fig:maj-pass}, reinforcement learning notably increases Maj@K accuracy but does not impact Pass@K. This observation suggests that reinforcement learning primarily strengthens the resilience of the output distribution, implying that \textbf{performance gains may largely stem from amplifying the likelihood of correct answers within the TopK candidates, rather than enhancing the model's core reasoning abilities.} This aligns with \citep{wang2023making}, who describe a \textbf{misalignment problem} encountered in SFT models on reasoning tasks, and demonstrate that aligning models' preferences—by employing various alignment strategies \citep{yuan2023rrhf,song2023preference,wang2023making}—can substantially improve their reasoning performance.

\subsubsection{How to Achieve More Effective RL?}
\label{sec:effec_rl}
Our results indicate that RL is highly effective for tasks involving mathematical reasoning. Furthermore, we introduce a unified framework that captures the essence of various prominent training strategies. In this framework, each approach can be interpreted as either a direct application or a streamlined variant of RL. As outlined in Equation \ref{eq:objective}, three principal elements structure this approach: the Data Source, the Algorithm, and the Reward Function. We suggest several possible research directions for improving each of these key components.

\paragraph{Data Source} The choice of data source serves as the foundational element for all training strategies. In the realm of RL, we define the data source as the collection of unlabeled questions paired with policy model-generated outputs. For the experiments described in this paper, our dataset comprises questions originating from the instruction tuning phase, and we employ a straightforward nucleus sampling method to produce outputs. We hypothesize that this limited approach may contribute to our RL pipeline yielding improvements primarily in Maj@K metrics. For subsequent work, we intend to apply our RL framework to prompts outside the training distribution and incorporate \textbf{advanced sampling (decoding) strategies} such as those leveraging tree-search methods \citep{yao2023tree}. Furthermore, advancements in \textbf{efficient inference techniques} \citep{xia-etal-2023-speculative,leviathan2023fast,kwon2023efficient,xia2024unlocking}, which are crucial for enhancing policy model exploration, are expected to play a pivotal role as well.

\paragraph{Algorithms} Algorithms utilize both input data and the reward signal to determine the gradient coefficient used in updating the model parameters. According to Equation \ref{eq:objective}, nearly all current approaches essentially \textbf{TRUST} the guidance provided by the reward function when adjusting the conditional probability of generating specific tokens. Nevertheless, the reliability of the reward signal cannot be consistently guaranteed, particularly in highly intricate tasks. For instance, even with the PRM800K datasets \citep{lightman2023let}—which are meticulously labeled by trained annotators—about 20\% of annotations are erroneous\footnote{\url{https://github.com/openai/prm800k/issues/12\#issuecomment-1728491852}}. Therefore, we focus on reinforcement learning algorithms that can withstand noisy reward signals. We posit that such \textbf{WEAK-TO-STRONG} \citep{burns2023weak} alignment strategies have the potential to significantly transform learning algorithms at their core.

\paragraph{Reward Function} The reward function serves as the key signal guiding the training process. In the context of RL, it is typically implemented as a neural reward model. We identify three pivotal areas for advancing reward models: 1) \textbf{Strategies to improve reward model generalization.} It is essential for the reward model to generalize well to unseen queries and novel decoding behaviors; failing to do so may result in reinforcement learning only consolidating the LLM’s output distribution without leading to substantive improvements in capability; 2) \textbf{Approaches for capturing and leveraging reward model uncertainty.} Incorporating uncertainty can potentially mediate between underperforming reward models and algorithms intended to enable progression from weaker to stronger learners; 3) \textbf{Practical methods for constructing precise, high-quality process reward models} that can deliver detailed and granular training feedback throughout the reasoning process \citep{lightman2023let,wang2023math}.

\section{Conclusion, Limitation, and Future Work} 

This work introduces DeepSeekMath, an open-source model that achieves superior results compared to all other open-source competitors on the MATH benchmark at the competition level, nearly matching the performance of proprietary systems. DeepSeekMath is based on DeepSeek-Coder-v1.5 7B and is further trained on an extensive dataset comprising 500B tokens, which includes 120B math-specific tokens primarily extracted from Common Crawl. Through comprehensive ablation studies, we determine that mathematical data harvested from web pages provides substantial advantages, whereas content from arXiv does not yield the anticipated improvements. We also propose Group Relative Policy Optimization (GRPO), an extension of Proximal Policy Optimization (PPO), capable of enhancing mathematical reasoning while using less memory. Experimental evidence demonstrates that GRPO offers effectiveness, even when DeepSeekMath-Instruct 7B attains strong benchmark results. Additionally, we present a unified framework for understanding related methodologies and identify several promising directions for further reinforcement learning research.

While DeepSeekMath~demonstrates strong performance on quantitative reasoning tasks, its proficiency in geometry and theorem proving falls short compared to proprietary models. For example, during our preliminary experiments, the model struggled with questions involving triangles and ellipses, suggesting possible biases in the data used for pre-training and fine-tuning. Moreover, due to the current model’s scale, DeepSeekMath~exhibits inferior few-shot learning abilities relative to GPT-4: whereas GPT-4 can leverage few-shot examples to enhance its results, DeepSeekMath~performs comparably in both zero-shot and few-shot settings. Moving forward, we plan to refine our engineered data selection process to assemble a higher-quality pre-training dataset. Furthermore, we will investigate strategies for more effective reinforcement learning of LLMs, as described in Section \ref{sec:effec_rl}.

\end{document}